\paragraph{参考评分与队列整理。}
本文队列由 \citet{xu2026sfibai} 所述机构超声数据库重新清洗整理而得，沿用专家给出的超声表现评分，范围为 0.0--3.5，间隔为 0.1。原研究的标注方案将四个临床等级细化为这一评分尺度，并说明了统一培训、集中复核和对评分分歧进行专家复核的流程。本研究对既有图像与标注进行清洗和整理，未重新评定评分。本文队列及划分与原研究不同，SFibAI 基线在本文划分上训练和评价。

\paragraph{扫查切面与病灶位置标注。}
六类切面标签从图像采集时记录的切面编号中提取，并在数据清洗时整理。病灶位置监督同样沿用 \citet{xu2026sfibai} 所述的既有医生矩形框，经清洗整理后使用。这些框标出图像中需要关注的区域，并附有原始分级标注。因此，切面标签描述采集时记录的扫查切面，矩形框则指示图像内部的局部区域。推理时无需提供这两类标注。

\paragraph{覆盖范围与弱目标。}
整理后的记录为每张图像保留切面及病灶位置标注，但并非每条记录都提供符合条件的非空弱目标。框监督使用目标分数为正的有效非空框，详见附录~\ref{sec:appendix-implementation}。弱定位评价包含 4,035 张有效测试图像，切面评价包含全部 4,107 张测试图像。局部框并未勾画完整病灶轮廓，因此框外响应的空间相关性仍需单独评价。
